% page 458 du fac-simile (image p-457 du scan)
% bloc 144.8 x 246.3 mm ; marge gauche 8.0 mm
\pgc{-12.700mm}{0.000mm}{
\l{\cel{3}450}
\l{}
\l{\sou{poliembriona}. (\sou{biol}.)Qua formacesis ek du o plura embri-}
\l{\cel{3}oni, da un ovulo. -}
\l{}
\l{\sou{polifonio}.\cel{2}(\sou{muziko)}. Efekto muzikala qua rezultas de la}
\l{\cel{3}ensemblo harmoniala di la instrumenti e di la voci.-}
\l{\cel{3}DEFIRS.}
\l{}
\l{\sou{poligama}.\cel{2}Spozulo qua havas plura spozini. DEFIRS.}
\l{}
\l{\sou{poliglota}. I. Redaktita ye plura lingui. - II. Qua savas}
\l{\cel{3}plura lingui. - DEFIRS.}
\l{}
\l{\sou{poligono.}\cel{2}I. (\sou{geom.)}\cel{2}Figuro qua havas plura anguli e}
\l{\cel{3}plura lateri. II. (\sou{termino militala)}; (a) figuro qua}
\l{\cel{3}determinas la trasuro di fortreso, di fortifikajo ;}
\l{\cel{3}(b) placo ube la artilriisti exercas su pri la pafo}
\l{\cel{3}per kanoni. - DEFIRS.}
\l{}
\l{\sou{poligrafo}. Autoro qua kompozas pri temi diversa. - DEFIS.}
\l{}
\l{\sou{polikroma}. Qua esas plura-kolora. - DEFIRS.}
\l{}
\l{\sou{polimatio}.\cel{2}Savo qua embracas plura cienci tre diversa.}
\l{\cel{3}DEFIS.}
\l{}
\l{\sou{polimorfa}. Qua chanjas ofte sua formo. - DEFIS.}
\l{}
\l{\sou{polinomio}. (\sou{mat.)}\cel{2}Quanto algebrala, ek plura termini qui}
\l{\cel{3}interligesas per la signi \textquotedbl{}plus\textquotedbl{} o \textquotedbl{}minus\textquotedbl{}. - DEFIS.}
\l{}
\l{\sou{polinuklea}. (\sou{histol.})\cel{2}Dicesas pri celulo o pri plastido}
\l{\cel{3}qua kontenas plura nuklei. -}
\l{}
\l{\sou{polipo}. I. (\sou{zool}.) Zoofito di qua la kapo havas periferie}
\l{\cel{3}palpili. - II. (\sou{patol.}) Surkreskajo qua su developas}
\l{\cel{3}en la kavaji qui havas membrano mukoza. - DEFIRS.}
\l{}
\l{\sou{poliporo}. (\sou{bot.}) Genero de fungi bazidiomiceta, karakt-}
\l{\cel{3}erizata da chapelo lignatra, qua havas sube tubi spo-}
\l{\cel{3}roza, paralela.}
\l{}
\l{\sou{polisar}. (\sou{trans}.)Igar glata per froto, netigar (armi,}
\l{\cel{3}koquilaro, e c.) per frotar per smerilo, greso, tri-}
\l{\cel{3}palio. - DEFIRS.}
\l{}
\l{\sou{polispermio}. (\sou{biol.}) Penetro en la ovulo matura di du o}
\l{\cel{3}plura spermatozoidi. -}
\l{}
\l{\sou{polita}. Karakterizata da lua egardo, indulgo, respekto}
\l{\cel{3}di sua kunhomi en mondo, en mondumo. - EFS.}
\l{}
\l{\sou{politeismo}. Religio qua agnoskas plura dei. - DEFIRS.}
\l{}
\l{\sou{politeisto}. La persono qua profesas politeismo. - DEFIRS.}
\l{politeknika. Qua koncernas plura arti, plura cienci.}
}
